\providecommand{\C}{\mathbb{C}}
\providecommand{\Q}{\mathbb{Q}}
\providecommand{\Z}{\mathbb{Z}}
\providecommand{\CoHA}{\mathrm{CoHA}}
\providecommand{\Hilb}{\mathrm{Hilb}}
\providecommand{\Perf}{\mathrm{Perf}}
\providecommand{\Rep}{\mathrm{Rep}}
\providecommand{\Sh}{\mathrm{Sh}}
\providecommand{\gglone}{\widehat{\mathfrak{gl}}_1}
\providecommand{\gghone}{\widehat{\widehat{\mathfrak{gl}}}_1}
\providecommand{\Wilpha}{\mathcal{W}_{1+\infty}}
\providecommand{\YFull}{Y}
\providecommand{\Yp}{Y^{+}}
\providecommand{\Ym}{Y^{-}}
\providecommand{\Yz}{Y^{0}}
\providecommand{\MO}{\mathrm{MO}}
\providecommand{\cZ}{\mathcal{Z}}

\section*{The positive CoHA of \texorpdfstring{$\C^3$}{C3} and the
\texorpdfstring{$\mathcal W_{1+\infty}$}{W1+infty} realisation}
\label{sec:coha-c3-w-infty-positive-half}

All statements are made over
\[
  \mathbb F
  =
  \Q(\epsilon_1,\epsilon_2,\epsilon_3)/
  (\epsilon_1+\epsilon_2+\epsilon_3).
\]
The quotient is the character field of the Calabi-Yau torus
\[
 T^{\mathrm{CY}}
 =
 \ker\bigl(\epsilon_1+\epsilon_2+\epsilon_3:
 (\C^\times)^3\to \C^\times\bigr)
\]
acting on \(\C^3=\operatorname{Spec}\C[x_1,x_2,x_3]\) with
\(\Omega_{\C^3}=dx_1\wedge dx_2\wedge dx_3\) fixed.  The third
weight is a record of the ambient three-coordinate action; on
\(T^{\mathrm{CY}}\) it is constrained by
\(\epsilon_3=-\epsilon_1-\epsilon_2\).

The objects which occur in the local \(\C^3\) calculation live in
different categories:
\[
\begin{array}{l|l|l}
\text{object} & \text{construction} & \text{structure} \\\hline
\CoHA_{T^{\mathrm{CY}}}(\C^3) & \text{critical CoHA of the Jordan triple loop quiver} &
E_1\text{-associative} \\
\Yp(\gglone) & \text{positive Feigin-Odesskii shuffle half} &
E_1\text{-associative} \\
\YFull(\gglone)=\Yp\bowtie\Yz\bowtie\Ym & \text{Hall-Drinfeld double} &
\text{topological Hopf algebra} \\
\mathcal E_{q_1,q_2,q_3}(\gglone) & \text{quantum toroidal algebra} &
\text{Hopf algebra with Miki }SL_2(\Z) \\
\Wilpha[\lambda] & \text{vertex algebra on the spectral curve} &
E_1\text{-chiral}
\end{array}
\]
The first two entries coincide.  The last entry is reached only after
doubling, central completion, an evaluation or Fock realisation, and
the state-field correspondence.  Thus
\[
  \CoHA(\C^3)=\Yp(\gglone),
  \qquad
  \CoHA(\C^3)\neq \Wilpha[\lambda],
  \qquad
  \Yp(\gglone)\neq \YFull(\gglone).
\]

\subsection*{The Calabi-Yau shuffle algebra}

Let \(\Lambda_n=\mathbb F[z_1,\ldots,z_n]^{S_n}\) and
\(\Lambda=\bigoplus_{n\geq0}\Lambda_n\).  The Feigin-Odesskii product is
\[
(F\star G)(x_1,\ldots,x_{m+n})
 =
\frac{1}{m!n!}
\sum_{\sigma\in S_{m+n}}
\sigma\!\left[
F(x_1,\ldots,x_m)G(x_{m+1},\ldots,x_{m+n})
\prod_{\substack{1\leq i\leq m\\m<j\leq m+n}}
\omega(x_j,x_i)
\right],
\]
where
\[
  \omega(z,w)
  =
  \frac{(z-w-\epsilon_1)(z-w-\epsilon_2)(z-w-\epsilon_3)}
       {(z-w)^3}.
\]
The kernel is written with three weights, but it is a two-parameter
Calabi-Yau kernel over \(\mathbb F\).  The self-dual identity
\(\epsilon_1+\epsilon_2+\epsilon_3=0\) is not a special point of the
family; it is the ambient Calabi-Yau constraint.

\begin{theorem}[Local positive-half theorem]
\label{thm:coha-c3-positive-half}
For the Jordan triple loop quiver
\[
Q=\bigl(\bullet \mathrel{\substack{\circlearrowleft X\\[-2pt]
\circlearrowleft Y\\[-2pt]\circlearrowleft Z}}\bigr),
\qquad
W=\operatorname{tr}(X[Y,Z]),
\]
the \(T^{\mathrm{CY}}\)-equivariant critical CoHA of \(\C^3\) is the
positive half of the affine Yangian of \(\widehat{\mathfrak{gl}}_1\):
\[
  \CoHA_{T^{\mathrm{CY}}}(\C^3)
  \cong
  \Sh^+_{\epsilon_1,\epsilon_2,\epsilon_3}
  \cong
  \Yp(\gglone)
\]
as \(E_1\)-associative algebras.  Under this normalisation the
Drinfeld current
\[
  E(u)=\sum_{k\geq0}E_k u^{-k-1}
\]
maps to the charge-one shuffle generators
\[
  E_k \longmapsto z^k\in \Lambda_1,\qquad k\geq0.
\]
\end{theorem}

\begin{proof}
Kontsevich and Soibelman define the critical CoHA from the vanishing
cycles of \(\operatorname{Tr}(W)\) on representation stacks of the
Jacobi algebra.  For the Jordan triple loop quiver the fixed-point
localisation on the \(T^{\mathrm{CY}}\)-equivariant Hilbert scheme
chart gives the shuffle product displayed above.  Schiffmann and
Vasserot identify the corresponding shuffle presentation on Hilbert
scheme cohomology with the positive part of the \(\mathfrak{gl}_1\)
affine Yangian, and Tsymbaliuk gives the Drinfeld-current
presentation with the same additive parameters.  Rapcak, Soibelman,
Yang, and Zhao formulate the same local toric CoHA as the positive
spherical CoHA whose double supplies the lowering operators.  These
identifications preserve the charge grading and send the charge-one
shuffle generator \(z^k\) to \(E_k\).
\end{proof}

The theorem is a positive-half statement.  The algebra
\(\bigoplus_n\Lambda_n\) contains the raising modes \(E_k\) and their
shuffle products.  It contains neither the lowering modes \(F_k\) nor
the Cartan tower \(\psi^\pm_k\) as internal elements.  A Chevalley
involution \(E_k\leftrightarrow F_k\) is defined only after passing to
the double; it does not restrict to an endomorphism of \(\Yp\).

\subsection*{The double and the vertex realisation}

The full current algebra is obtained by a Hall-Drinfeld double with
respect to a nondegenerate Euler or stable-envelope pairing:
\[
  \mathcal D_\hbar(\Yp)
  =
  \Yp\bowtie\Yz\bowtie\Ym
  =
  \YFull(\gglone).
\]
In Drinfeld currents,
\[
 E(u)=\sum_{k\geq0}E_k u^{-k-1},\qquad
 F(u)=\sum_{k\geq0}F_k u^{-k-1},\qquad
 \psi^\pm(u)=1+\sum_{k\geq0}\psi^\pm_k u^{\mp k\mp1},
\]
with structure function
\[
  g(u)
  =
  \frac{(u-\epsilon_1)(u-\epsilon_2)(u-\epsilon_3)}
       {(u+\epsilon_1)(u+\epsilon_2)(u+\epsilon_3)}.
\]
The Calabi-Yau condition gives \(g(u)g(-u)=1\), the unitarity
normalisation of the rational \(R\)-matrix.  Expanding
\(\log g(u)\) at \(u=\infty\) leaves only odd powers at
\(\epsilon_1+\epsilon_2+\epsilon_3=0\), since the even power sums cancel
in the numerator-denominator ratio.

The passage to \(\Wilpha[\lambda]\) is not another name for
\(\CoHA(\C^3)\).  It is the completed chain
\[
\CoHA(\C^3)
\cong
\Yp(\gglone)
\hookrightarrow
\mathcal D_\hbar(\Yp)
=
\YFull(\gglone)
\rightsquigarrow
\widehat{\YFull}_{\mathrm{eval}}
\cong
\operatorname{Modes}(\Wilpha[\lambda]).
\]
The last arrow means central completion plus an evaluation or Fock
module on which the state-field map is defined.  In the standard
normalisation the raising and lowering modes correspond to opposite
Fourier halves of the vertex-algebra fields:
\[
  E_k \leftrightarrow J^{(2)}_{-k-1},
  \qquad
  F_k \leftrightarrow J^{(2)}_{k+1},
\]
while \(\psi^\pm\) record the Heisenberg-Cartan tower.  A vertex
algebra has fields
\[
  J^{(s)}(z)=\sum_{n\in\Z}J^{(s)}_n z^{-n-s},
  \qquad s\geq1,
\]
and therefore requires both Fourier halves.  The positive CoHA supplies
only the creation half.

The \(d=3\) output of the Calabi-Yau-to-chiral construction is
\(E_1\)-chiral:
\[
  A_{\C^3}=\Phi_3^{(\Sigma_2,C)}(\Perf(\C^3))
  \in E_1\text{-}\mathrm{ChirAlg}.
\]
On constructed loci the \(E_2\)-braided structure lives on
\[
  \cZ\bigl(\Rep^{E_1}(A_{\C^3})\bigr),
\]
not on \(A_{\C^3}\) itself and not on the positive CoHA\@.  This is the
operadic boundary that prevents
\(\CoHA(\C^3)=\Wilpha[\lambda]\).

\subsection*{Stable envelopes}

For \(\Hilb^n(\C^2)\), Maulik-Okounkov stable envelopes give chamber
maps
\[
  \operatorname{Stab}_{\mathfrak C}:
  H^*_T\bigl(\Hilb^n(\C^2)^T\bigr)\to H^*_T(\Hilb^n(\C^2)).
\]
The ratio
\[
  R_{\MO}(u)=\operatorname{Stab}_{-}^{-1}\operatorname{Stab}_{+}
\]
satisfies the Yang-Baxter equation and has the scalar rank-one
structure function \(g(u)\) displayed above.  Its full fixed-point
matrix is triangular with arm-leg weight factors; it is not a diagonal
product formula in general.  Increasing Nakajima correspondences give
the \(E_k\)-half, decreasing correspondences give the \(F_k\)-half, and
the stable-envelope pairing is the geometric source of the double:
\[
  \langle \alpha,\beta\rangle_{\MO}
  =
  \int_{\Hilb^n(\C^2)}^T
  \operatorname{Stab}_{+}(\alpha)\cup
  \operatorname{Stab}_{-}(\beta).
\]
Thus stable envelopes restore the missing half geometrically; they do
not turn \(\Yp\) itself into a vertex algebra.

\subsection*{Parameters and the Heisenberg fibre}

The equation \(\epsilon_1+\epsilon_2+\epsilon_3=0\) leaves a projective
one-parameter family of Calabi-Yau weights.  It does not force the
\(\Wilpha\)-parameter to a single value.  The special point
\[
  \epsilon_3=0,\qquad \epsilon_1=-\epsilon_2
\]
is the Heisenberg fibre.  There the shuffle kernel becomes
\[
  \omega(z,w)
  =
  \frac{(z-w-\epsilon_1)(z-w+\epsilon_1)(z-w)}
       {(z-w)^3}
  =
  \frac{(z-w)^2-\epsilon_1^2}{(z-w)^2}.
\]
The positive half generated by \(z^k\in\Lambda_1\) is the creation
half of the level-one Heisenberg algebra; after doubling and
state-field realisation one obtains the Heisenberg VOA\@.  The local
normalisation
\[
  \kappa_{\mathrm{ch}}^{\mathrm{Heis}}(\C^3)=1
\]
is the Heisenberg normalisation for the local \(d=3\) chart.
It is not \(c/2\), not a compact Hodge supertrace of \(\C^3\), and not
a Borcherds weight.

\subsection*{Conifold and local \texorpdfstring{$\mathbb P^2$}{P2}}

The resolved conifold is a different toric CY\(_3\).  Its positive CoHA
has the Negut shuffle kernel
\[
  \omega_{\mathrm{con}}(z,w)
  =
  \frac{(z-w+\epsilon_1)(z-w+\epsilon_2)}
       {(z-w)(z-w+\epsilon_1+\epsilon_2)}.
\]
Davison's conifold CoHA and the Rapcak, Soibelman, Yang, and Zhao
toric construction identify this positive half with a conifold
Yangian or quantum toroidal \(\widehat{\widehat{\mathfrak{gl}}}_2\)
object, after the appropriate stability and shift data are fixed.  It
is not \(\Yp(\gglone)\) and not \(\Wilpha(\gglone)\).

Perturbative five-dimensional Chern-Simons diagrams, shuffle
multiplication, and vertex OPEs agree only after the relevant
doubling, completion, and state-field realisation have been specified.
At the \(\C^3\) Heisenberg fibre this produces the usual
\(\Wilpha\)-Heisenberg realisation.  At the conifold it produces a
conifold-specific quantum-toroidal vertex algebra.

For local \(\mathbb P^2\), the positive half is controlled by the
\(\Z_3\)-McKay quiver.  The relevant Yangian is a quiver Yangian
attached to that McKay data, not the \(\widehat{\mathfrak{gl}}_1\)
positive half of \(\C^3\).  Orbifold or coset descriptions of the
vertex algebra occur only after the same double-plus-state-field
passage.

\subsection*{\texorpdfstring{$K3\times E$}{K3 x E}}

The compact threefold \(K3\times E\) is not a toric local chart.  The
\(\C^3\) fixed-point shuffle proof does not globalise to it.  The
correct compact source is a reduced critical Hall object built from
K3-cosection-reduced moduli and finite Hall windows.  The local
Davison-Meinhardt PBW theorem says that, under the quiver-with-potential
and critical-chart hypotheses, a critical CoHA has a PBW degeneration
whose primitive part is a BPS Lie algebra:
\[
  \operatorname{gr}\CoHA
  \cong
  \operatorname{Sym}\bigl(\mathfrak g_{\mathrm{BPS}}\otimes H^*(B\C^\times)\bigr).
\]
This is PBW control by the BPS primitive part, not an identification
with a prescribed BKM triangular half.
For compact \(K3\times E\), finite-window Hall objects must still pass
orientation, parity, radical, PBW, centre, and completion gates before
a Hall-Borcherds double is recognised.  The
unconditional numerical evidence is character-level; the bracket-level
comparison
\[
  \mathfrak g_{\mathrm{BPS}}(K3\times E)
  \stackrel{?}{\cong}
  \mathfrak g_{\Delta_5}
\]
remains an open recognition problem outside finite windows where all
defects have been killed.

The automorphic scalar is separate from the Hall multiplication.  The
Gritsenko-Nikulin Borcherds product gives
\[
  \Phi_{10}^{\mathrm{un}}=\Delta_5^2,
  \qquad
  c_1(0)=10,
  \qquad
  \kappa_{\mathrm{BKM}}(\Delta_5)=\frac{c_1(0)}2=5.
\]
In Oberdieck-Pandharipande's reduced DT normalisation one writes
\[
  D_5=64^{-1}\Delta_5,
  \qquad
  \Phi_{10}^{\mathrm{OP}}=D_5^2=4096^{-1}\Delta_5^2,
  \qquad
  Z^{\mathrm{red}}_{\mathrm{DT}}(K3\times E)
  =-\bigl(\Phi_{10}^{\mathrm{OP}}\bigr)^{-1}.
\]
These identities determine denominator and character data.  They do
not imply
\[
  \CoHA(K3\times E)
  =
  U\bigl(\mathfrak n_+(\mathfrak g_{\Delta_5})\bigr).
\]
The BKM-side compact object is a Hall-Drinfeld double or finite
radical-quotient double of the reduced compact Hall source, when the
recognition data exist.  It is not the local \(\C^3\) affine Yangian
and not the K3 self-mirror Yangian branch.

\begin{theorem}[Taxonomy]
\label{thm:coha-w-taxonomy}
The following identifications and non-identifications are forced by
the constructions above.
\[
\CoHA_{T^{\mathrm{CY}}}(\C^3)
\cong
\Yp(\gglone),
\qquad
\mathcal D_\hbar(\Yp)
=
\YFull(\gglone),
\]
\[
\widehat{\YFull}_{\mathrm{eval}}
\cong
\operatorname{Modes}(\Wilpha[\lambda]),
\qquad
\CoHA(\C^3)\neq \Wilpha[\lambda].
\]
For \(d=3\), \(\Phi_3\) produces \(E_1\)-chiral algebras after the
chosen specialisation; the \(E_2\)-braided structure is a centre or
representation-category structure.  For \(K3\times E\), the
\(\Delta_5\) denominator fixes \(\kappa_{\mathrm{BKM}}=5\) and
character data, while the compact Hall-Borcherds algebra-level
identification remains conditional on the finite recognition data.
\end{theorem}

\begin{proof}
The first isomorphism is Theorem~\ref{thm:coha-c3-positive-half}.  The
second is the Hall-Drinfeld double construction: it adds the missing
lowering half and Cartan tower through the Euler or stable-envelope
pairing.  The completed evaluated double acts as the affine Yangian
mode algebra of \(\Wilpha[\lambda]\), where the state-field
correspondence supplies the vertex operation.  Since a vertex algebra
has all Fourier modes and locality, while \(\CoHA(\C^3)\) has only the
positive shuffle product, the equality \(\CoHA(\C^3)=\Wilpha[\lambda]\)
is category-theoretically impossible.  The dimension-three operadic
scope follows from the construction of \(\Phi_3\): the native object is
\(E_1\)-chiral and its \(E_2\) enhancement, when present, is carried by
the Drinfeld centre.  The \(K3\times E\) statement follows from the
separation between Borcherds denominator arithmetic and compact Hall
multiplication.
\end{proof}

\subsection*{Excluded equalities}

The corrected local and compact formulas are
\[
  \CoHA(\C^3)=\Yp(\gglone),
  \qquad
  \mathcal D_\hbar(\Yp)=\YFull(\gglone),
  \qquad
  \widehat{\YFull}_{\mathrm{eval}}
  \cong
  \operatorname{Modes}(\Wilpha[\lambda]),
\]
\[
  \kappa_{\mathrm{BKM}}(\Delta_5)=5,
  \qquad
  \Phi_{10}^{\mathrm{un}}=\Delta_5^2.
\]
The false shortcuts are
\[
  \CoHA(\C^3)=\Wilpha[\lambda],
  \qquad
  \Yp(\gglone)=\YFull(\gglone),
  \qquad
  \CoHA(K3\times E)
  =
  U\bigl(\mathfrak n_+(\mathfrak g_{\Delta_5})\bigr)
  \text{ unconditionally}.
\]
The local \(\C^3\) positive half, the compact \(K3\times E\)
Hall-Borcherds branch, and the K3 self-mirror Yangian branch are
different constructions.  Their characters can meet; their algebra
structures do not become equal without the explicit double,
completion, centre, and recognition data.
